\clearpage
\section{E4: Stage 2 특징이 실제 팩 가치로 이어지는가?}
\state{2026년 11월 5일 설계. 입력 점검·실제 용량 대조군 실행 완료. 추가 특징 모델 비교는 미실행.}

\subsection*{질문과 비교 조건}
팀 피드백의 질문은 쿨롱 효율, dQ/dV 곡선, 분극·저항 등의 특징으로 줄어든 수명 예측 오차가 실제 팩 가치 $V$ 개선으로 이어지는가이다. 먼저 기본 진단 10사이클에서 비교하고, 이후 7개 공통 예산으로 확장한다. 86셀 전체·팩 크기·진단 비용·미사용 규칙은 E3를 그대로 따른다.

\begin{center}
\begin{tabular}{p{.21\linewidth}p{.68\linewidth}}
\toprule
방법 & 입력과 비교 목적 \\
\midrule
Current-only & 전 셀 $n_i=10$, 사이클 6 용량 정렬. 비용 860의 고정 기준선. \\
용량 특징 모델 & 기존 E2-B의 용량 특징 10개를 2개 RPT에서도 계산하도록 명시적으로 수정한 V3 모델. \\
Stage 2 모델 & 팀의 screened 특징 20개. 용량 모델과 같은 모델 후보·분할·튜닝·배정 절차로 입력 효과를 비교. \\
최신 용량 정렬 & 보조 대조군. 각 균등 진단량에서 마지막 공개 용량으로 정렬. 추가 진단 비용까지 같은 조건에서 비교. \\
\bottomrule
\end{tabular}
\end{center}

두 학습 방법은 Ridge와 GP 각각의 결과를 나란히 보고하고, 동일한 내부 검증 기준으로 선택한 모델의 결과도 별도로 보고한다. 서로 다른 모델을 선택한 두 결과만으로 특징의 효과를 단정하지 않는다. 첫 배정 규칙은 예측 RUL 내림차순 정렬 후 연속 묶기이며, 하위 나머지 셀은 미사용으로 남긴다. 추가 특징이 예측한 배정만 바꾸고 평가는 실제 미래 용량으로 계산하므로, \textbf{이 1차 비교에는 미래 용량 궤적 예측모델이 필수는 아니다.}

\subsection*{관측 시점과 목표값}
$n_i=n$이면 공개 RPT는 $1,6,\ldots,n-4$이고 재사용 평가 시작은 $n+1$이다. 모델 목표는 이 시작 시점부터의 $\mathrm{RUL}=T_{\mathrm{eol}}-(n+1)$로 통일한다. 기존 E2-B가 마지막 관측 사이클부터 계산했던 RUL과는 기준점이 5사이클 다르다. 예측 점은 두 모델 모두 로그 RUL 예측값의 지수, 즉 중앙값으로 통일한다.

기존 용량 특징의 초기·최근 3점 기울기는 관측이 두 개이면 두 점 기울기로 대체한다. 이때 전체·초기·최근 기울기는 같고 기울기 변화와 직선 잔차는 0이다. 이를 기존 E2-B의 재현으로 부르지 않고 V3 입력 경계 수정으로 기록한다.

\subsection*{공유 fold와 평가 범위}
공유 셀 단위 5-fold를 사용한다. 외부 평가 fold를 제외하고, 내부 3-fold에서 결측 처리·표준화·모델 선택을 수행한다. 시드 11,29,47은 두 방법에 공통으로 적용한다. 설계 시점에는 시드 42 임시본을 사용했으나, 최신 `origin/main`에는 별도 공유 fold가 있다. 과거 임시본 기준 결과와 새 fold 기준 결과를 섞지 않는다.

전체 86셀에 대해 각 셀을 학습에서 제외한 OOF 예측을 모으면 86셀 팩 배정을 탐색할 수 있다. 다만 서로 다른 셀의 예측에 다른 fold의 정답이 학습으로 들어가므로, 이를 독립된 86셀 평가 재고의 일반화 성능으로 주장하지 않는다. OOF 전체 재고 결과는 탐색용, 각 평가 fold의 예측 오차는 셀 분리 평가로 표시한다. fold 내 팩 가치는 17~18셀의 보조 분석으로 따로 보고하며 $N=86$ 본 실험과 합치지 않는다. 완전 독립 86셀 팩 평가에는 별도 학습 셀 또는 별도 평가 재고가 필요하다.

\clearpage
\subsection*{입력 준비와 진행 기준}
설계 당시에는 \file{features_v2.npz}와 \file{final_selection.json}이 없어 실제 키·축을 확인하지 못했다. 최신 `origin/main`에는 두 파일이 추가되어 있다. 현재 확인한 배열 키는 \texttt{X}, \texttt{cell\_ids}, \texttt{budgets}, \texttt{feature\_names}, \texttt{base\_names}이며 $X$의 크기는 $86\times7\times148$이고 셀 ID는 \file{C\_full.npz}와 일치한다. 선정 JSON은 \texttt{screened} 20개와 별도 \texttt{features} 6개를 담는다. V3 설정에서 진단량 키를 \texttt{budgets}에 매핑했다.

입력 점검기는 셀 ID 정렬, 진단량 축 10~40, 배열 크기, 특징 이름 중복, 20개 선택, 파일 해시와 결측 비율을 확인하도록 준비했다. 실제 `budgets` 키로 preflight를 다시 실행해야 하며, 그 결과가 나오기 전까지 입력 준비 완료를 주장하지 않는다. 행 번호를 셀 ID로 간주하거나 배열 축을 임의 추정하지 않는다. 진단량별 특징은 해당 시점 공개 관측만으로 계산되어야 한다. 148개 중 20개를 선정할 때 평가 셀의 정답을 사용했다면 외부 fold의 학습 셀 안에서 다시 선정해야 한다. 현재 제공된 20개를 고정해 탐색하는 결과는 독립 평가와 구분한다.

\begin{enumerate}
\item 입력: 특징 파일·선정 목록 수신, 실제 형식 매핑, 공개 관측 범위와 선정 근거 확인.
\item 예측: 동일 fold·내부 검증·모델 후보로 두 특징군의 OOF 예측과 fold별 오차를 저장.
\item 배정: 같은 정렬 규칙으로 4·8·12·16직렬 팩과 미사용 셀을 정하고 실제 $V$ 평가.
\item 확장: 7개 균등 예산 곡선을 만든 뒤, 예측모델을 고정해 검사 선택 전략만 비교.
\end{enumerate}

현재 준비된 실행 코드는 1단계 입력 점검과 실제 용량 대조군이다. 학습·튜닝 실행기는 아직 연결하지 않았으며 특징 파일을 추가하는 것만으로 전체 실험이 자동 실행되는 상태는 아니다.

\subsection*{판정 지표와 oracle 처리}
주 비교는 $n=10$에서 Stage 2와 Current-only, Stage 2와 용량 특징 모델 간 $\Delta V$이다. 네 팩 크기별 절대 $V$, $\Delta V$, 총 진단 비용을 모두 보고한다. 예측 MAE가 줄어도 $V$가 늘지 않으면 예측 개선이 배정 이득으로 이어졌다고 주장하지 않는다. 같은 셀·같은 시드의 짝 비교를 유지하되, 시드 세 개를 독립 재고 표본처럼 취급하지 않는다.

총 regret은 E3의 $V_O$가 확정된 뒤 보고한다. 당장은 $V$와 차이를 계산하고, 미정 regret을 0으로 채우지 않는다. 향후 다중 시작 교환 탐색은 실행 가능한 배정의 하한, 시점별 재그룹화는 완화 상한으로 구분해 저장하고 전역 최적 여부와 간격을 명시해야 한다. 과거 E2-B 점수는 구지표 부록으로만 비교한다.

\subsection*{후속 적응 진단 설계}
예측기와 팩 구성 절차를 고정하고 균등·무작위·RUL 불확실성·DFF 선택만 바꾼다. 매 라운드는 현재 관측에서 예측, 시나리오 배정, DFF 계산, 9셀 선택, 5사이클 진단 비용 반영, 새 관측 공개의 순서이다. 미사용 후보도 선택 가능하며 셀당 40사이클 상한과 공통 예산을 지킨다. 라운드 내 선택은 공개 전 한 번에 결정한다.

DFF의 기준 배정은 모델 예측으로 정한다. 참 미래 oracle을 선택 정책에 전달하지 않는다. DFF는 동료 변동성에 기반한 검사 선택 후보이며 보상 자체가 아니다. 진단 후 실제 $V$가 개선되는지로 검증한다. RUL 시나리오의 최소값 합을 최적화하는 교환은 $V$ 최적화와 구분한다. 예측 $V$ 기반 교환·EVOI를 주장하려면 미래 용량 궤적과 그 불확실성 모델을 추가로 설계해야 한다. 중단 규칙은 고정 예산 비교 후 검토한다.

\clearpage
\subsection*{실제로 실행한 대조군과 검증}
실제 \file{C_full.npz}에서 균등 진단 후 최신 공개 용량으로 정렬한 결과이다. 아래 값은 예측모델이나 DFF 성능 결과가 아니다. $n=10$의 최신 용량 정렬은 Current-only와 같다. Current-only는 이후 예산에서도 비용 860의 고정 비교점으로 유지하며, 추가 진단을 한 것처럼 비용을 바꾸지 않는다.

\begin{center}
\begin{tabular}{rrrrrr}
\toprule
셀당 진단 & 총 비용 & 4직렬 $V$ & 8직렬 $V$ & 12직렬 $V$ & 16직렬 $V$ \\
\midrule
10 & 860   & 7704.393 & 3374.804 & 2104.898 & 1477.988 \\
15 & 1290  & 7369.421 & 3139.121 & 2067.727 & 1463.776 \\
20 & 1720  & 7605.423 & 3042.367 & 1963.880 & 1392.185 \\
25 & 2150  & 6941.466 & 3120.218 & 2051.037 & 1393.471 \\
30 & 2580  & 6761.720 & 3012.233 & 1975.626 & 1348.915 \\
35 & 3010  & 6769.517 & 2865.628 & 1857.594 & 1303.974 \\
40 & 3440  & 6476.763 & 2826.378 & 1793.846 & 1276.047 \\
\bottomrule
\end{tabular}
\end{center}

$V$의 단위는 Ah·사이클, 비용 단위는 셀-사이클이다. 진단이 늘면 용량 정렬 정보와 재사용 시작점이 동시에 바뀐다. 따라서 이 곡선만으로 진단 손실과 배정 손실을 분리하지 않는다. 추가 진단과 재정렬이 자동으로 $V$를 높이지 않는다는 대조군 결과로 해석한다.

실행은 28개 최신 용량 정렬과 4개 Current-only, 총 32개 조건을 저장했다. 원시 결과와 배정 ID는 \file{results/experiment_design_v3/control_metrics.csv}, \file{control_assignments.json}에 있다. 파일 존재·해시·fold·준비 상태는 \file{readiness.json}에 기록한다. 기존 검증을 포함한 35개 검증이 통과했으며, 새 검증은 두 RPT 특징 계산, 미래값 변경에 대한 입력 불변성, 셀 ID 재정렬, Current-only와 초기 대조군 일치를 확인했다.

\subsection*{VS Code에서 재현하는 순서}
\begin{enumerate}
\item \file{simulation} 폴더를 열고 실행 메뉴의 \textbf{V3. 새 실험 입력·준비 상태 확인}을 선택한다. 기본 빌드 작업도 이 점검으로 변경했다.
\item \textbf{V3. 실제 용량 기준선 7개 예산 실행}을 선택하면 현재 데이터로 위 표와 배정 ID를 다시 저장한다.
\item \file{configs/experiment_plan_v3.json}에서 데이터·특징·선정 파일 경로, 진단량·팩 크기, 모델 후보와 시드를 검토한다. 현재 완성된 실행과 후속 설계 항목을 구분해 표시했다.
\item 특징 파일 수신 후 실제 키·특징 목록과 선정 근거를 확인하고 학습 실행기를 연결한다. V1/V2 실행 메뉴는 과거 재현용으로 보존한다.
\end{enumerate}

터미널에서는 \file{simulation}을 기준으로 다음 두 명령을 실행한다. 긴 구현 코드는 이 문서에 넣지 않았다.
\begin{Verbatim}[fontsize=\small]
uv run --directory simulation python -m battery_sim.experiment_v3 --stage preflight
uv run --directory simulation python -m battery_sim.experiment_v3 --stage controls
\end{Verbatim}
